\documentclass[10pt]{article}

\usepackage[utf8]{inputenc}
\usepackage[export]{adjustbox}
\usepackage{fancyhdr}
\usepackage{hyperref}
\usepackage{datetime}
\usepackage{natbib}
\usepackage{array}
\usepackage{amsmath,amssymb}
\usepackage[includeheadfoot,
            left=0.85in,
            right=0.85in,
            top=-0.55in,
            bottom=0.05in,
            headheight=110pt]{geometry}
\usepackage{titlesec}

\newdateformat{monthyeardate}{%
  \monthname[\THEMONTH] \THEYEAR}

\pagestyle{fancy}
\fancyhf{}
\renewcommand{\headrulewidth}{0pt}
\renewcommand{\footrulewidth}{0pt}
\renewcommand{\refname}{}
\fancyhead[L]{
    \begin{minipage}[t]{0.55\textwidth}
    \vspace{0pt}
    \includegraphics[height=28pt]{logo.pdf}
    \end{minipage}
    \hfill
    \begin{minipage}[t]{0.4\textwidth}
    \raggedleft
    \vspace{0em}
    Alejandro Trigueros Alonso \\
    Felicia Scharitzer
    \end{minipage}
    \vspace{-0.75\baselineskip}
}

\setlength{\footskip}{20pt}
\fancyfoot[R]{\monthyeardate\today}

\renewcommand{\thesection}{Part \arabic{section}}
\renewcommand{\thesubsection}{\arabic{subsection} .}

\begin{document}

\setlength{\headheight}{8\baselineskip}
\setlength{\parindent}{0pt}
\setlength{\parskip}{\baselineskip}

\titlespacing*{\section}{0pt}{0.0ex}{0.0ex}
\titlespacing*{\subsection}{0pt}{0ex}{-1ex}

\section*{Conformal Prediction under High-Dimensional Covariate Shift}
\vspace{-1em}
{\itshape Guarantees in Machine Learning, Project Proposal. Supervisor: Francesco Taroni}

\subsection*{Main paper and field}

\vspace{-0.25em}

This project studies conformal prediction under covariate shift: how to keep prediction
sets reliable when the model is deployed on inputs that differ from its calibration data.
The main paper is \textit{Conformal Inference under High-Dimensional Covariate Shifts via
Likelihood-Ratio Regularization} by Joshi et al.\ 2025 \cite{joshi2026conformal}.

\subsection*{Motivation}

\vspace{-0.25em}

Conformal prediction turns any trained model's output into prediction sets that contain
the true value with probability at least $1{-}\alpha$. Split conformal prediction, the
most common variant, computes the model's errors on a held-out calibration set and takes
their $(1{-}\alpha)$-quantile as the set width. The guarantee holds in finite samples and
without assumptions on the model or the data distribution, requiring only that calibration
and test points are exchangeable.

\vspace{-0.25em}

This assumption fails under covariate shift: the input distribution changes from
$\mathbb{P}_1$ (calibration) to $\mathbb{P}_2$ (test) while the conditional $P(Y \mid X)$
remains the same. The model can be far less accurate on the test inputs, and coverage
drops below the nominal level without warning. We are interested in the setting where a
neural network surrogate replaces an expensive simulation inside a design optimizer. The
conditional relationship does not change, but the optimizer explores inputs that differ
from the training data, often in regions where the surrogate is least reliable, producing
exactly the covariate shift setting with labeled calibration data from $\mathbb{P}_1$ and
only unlabeled inputs from $\mathbb{P}_2$.

\vspace{-0.25em}

Three families of methods address this setting. The first reweights calibration errors by
the likelihood ratio $r = d\mathbb{P}_2/d\mathbb{P}_1$, restoring exact coverage when $r$
is known \cite{tibshirani2019conformal}. When $r$ must be estimated, coverage degrades
with the estimation error \cite{lei2021conformal, yang2024doubly}. An alternative fits a
quantile regression over a function class, giving coverage under every shift in that class
simultaneously without requiring target data \cite{gibbs2025conformal}, at the cost of
over-coverage of order $\mathcal{O}(d/(n{+}1))$. A third family, distributionally robust
methods, guarantees coverage under the worst-case shift within a neighborhood of the
calibration distribution \cite{aolaritei2026conformal}. When inputs are high-dimensional
(images, embeddings, physical geometries) and labeled data is expensive, all three face
difficulties: density ratio estimation becomes unreliable, class-wide coverage is
conservative, and robust neighborhoods require labeled test data to calibrate.

\subsection*{Paper's methodology}

\vspace{-0.25em}

Joshi et al.\ address the case where inputs are high-dimensional and estimating $r$
directly is impractical. Their method (LR-QR) fits an input-dependent threshold $h(x)$
through a set of features $\Phi(x)$ (for example the embedding of a pretrained network)
by minimizing
\[
\min_{h\in\mathcal H,\,\beta\in\mathbb R}\ \mathbb E_1[\ell_\alpha(h(X),S)]
+ \lambda\big(\beta^2\,\mathbb E_1[h(X)^2]-2\beta\,\mathbb E_2[h(X)]\big),
\]
where $S$ is the conformity score of a labeled calibration point and $\ell_\alpha$ is the
pinball loss, whose minimizer is the $(1{-}\alpha)$-quantile. The first term is a quantile
regression on calibration data, as in \cite{gibbs2025conformal}, but instead of
guaranteeing coverage under all shifts in the function class, the second term uses
unlabeled inputs from $\mathbb{P}_2$ to focus on the specific shift that occurs. Since
$\mathbb E_2[h]=\mathbb E_1[r\,h]$, this term equals the squared distance between
$\beta h$ and $r$ up to a constant, steering the threshold towards the actual shift and
avoiding the over-coverage penalty of the class-wide guarantee. At the population level,
coverage at $1{-}\alpha$ holds assuming $r$ lies in the span of the features, with
additional estimation error and misspecification terms in finite samples. The paper
validates the method on tabular data, an LLM question-answering task, and cell images
with synthetic shifts, but does not consider physical simulation data.

\vspace{-0.25em}

The goal of this project is to clarify the relationship between LR-QR and the results it
builds on, and to assess whether the method's assumptions and practical choices (feature
selection, regularization strength, shift magnitude) transfer to application settings like
surrogate-based design optimization, where the shift is induced by an optimizer and inputs
are high-dimensional.

\subsection*{Related Work}

\vspace{-0.25em}

\textit{Conformal prediction under covariate shift} \cite{tibshirani2019conformal}.
Introduces the reweighting approach and proves exact coverage when the likelihood ratio is
known. This is the starting point for the problem LR-QR addresses and one of its
baselines.

\vspace{-0.25em}

\textit{Conformal prediction beyond exchangeability} \cite{barber2023conformal}. Extends
conformal prediction to non-exchangeable data, bounding the coverage loss by how far the
data depart from exchangeability. It provides the general framework for analysing methods
under shift.

\vspace{-0.25em}

\textit{Conformal prediction with conditional guarantees} \cite{gibbs2025conformal}.
Shows that quantile regression over a $d$-dimensional function class guarantees coverage
under every shift in that class simultaneously, without target data, at the price of sets
that over-cover by up to order $d/(n+1)$. LR-QR builds on this construction and adds
unlabeled target inputs to focus on the shift that actually occurs.

\vspace{-0.25em}

\textit{A least-squares approach to direct importance estimation} \cite{kanamori2009least}.
Estimates the likelihood ratio by minimizing $\tfrac12\mathbb E_1[g^2]-\mathbb E_2[g]$,
which up to a factor of two is the LR-QR regularizer with $g=\beta h$. The authors of
LR-QR show that their fitted $\beta h$ stays far from $r$; we want to understand what the
regularizer actually does and how much of the method follows from these existing pieces.

\subsection*{Next Steps}

\vspace{-0.25em}

\begin{enumerate}
    \item Position LR-QR in the literature: what problem it addresses that earlier methods
          left open, and how later work relates to it.
    \item Work through the paper's proofs, assumptions, and experiments, with LLM- and
          Lean-assisted checking.
    \item Identify which assumptions and practical choices might not transfer to
          surrogate-based design: whether the features capture the shift, how the
          regularization strength is chosen, and how optimization-induced shifts are
          handled.
    \item Select a concrete question from the above and develop it, combining theoretical
          analysis with small-scale experiments.
\end{enumerate}

\vspace{-5ex}

\bibliographystyle{unsrt}

\bibliography{references}

\end{document}
